% GENERATED FILE. Edit canonical lessons, then run npm run generate:books.
% canonical-source-hash: fnv1a64:72fe3742c8773e05
% canonical-lessons: ZH-C537-qiyin, ZH-C537-yuanlai, ZH-W537-yuan-origin, ZH-C537-daozhi, ZH-C537-zaocheng, ZH-C537-xiaoguo

\chapter{As It Turns Out}
\label{ch:zh-as-it-turns-out}
\hlchaptermodality{537}

\begin{chapteropening}
\textbf{By the end of this chapter:} \emph{I can say what started something, that it turns out to be so, what it leads to and causes, and whether it works.}

Complete the last lesson of chapter 537: I can say what started something, that it turns out to be so, what it leads to and causes, and whether it works.
\end{chapteropening}

\section[\texorpdfstring{\zh{起因} (qǐyīn)}{起因 (qǐyīn)}]{\zh{起因} (qǐyīn) — what set it off; the starting cause}

\label{lesson:ZH-C537-qiyin}



\begin{quote}
Before the new one: say the Mandarin for to see in a dream, then the Mandarin for something to look forward to.
\end{quote}

\subsection*{You'll want to know: \zh{起因}}
\textbf{\zh{起因}} — \emph{qǐyīn} — \textquotedblleft{}what set it off; the starting cause\textquotedblright{}.

\begin{quote}
\emph{rúguǒ zhīdào le qǐyīn jiù néng jiějué wèntí} — If we know what set it off, we can solve the problem.
\end{quote}

\begin{cousinweb}
\textbf{\zh{起因}} (\emph{qǐyīn}) is built from \zh{起} 'to rise, to start' and \zh{因} 'a cause', both inherited from Middle Chinese, \emph{khiX} and \emph{'jin} (Guangyun, in Baxter's transcription), which glosses the first 'to rise' and 'to start' and the second 'to rest on' and 'to follow from'.
\end{cousinweb}

\begin{grammarlens}[title={What you've built}]
One more word for saying what would happen.
\end{grammarlens}

\subsection*{Your turn}
\begin{itemize}
\raggedright
  \item \emph{Say it:} \emph{qǐyīn}
  \item \emph{Say it:} \emph{qǐyīn}, once more
  \item \emph{Say it:} \emph{rúguǒ zhīdào le qǐyīn jiù néng jiějué wèntí} — If we know what set it off, we can solve the problem.
  \item \emph{Recall:} say the Mandarin for to see in a dream, then the Mandarin for something to look forward to, then say \emph{qǐyīn} again
\end{itemize}

\begin{tcolorbox}[breakable,colback=teal!4,colframe=teal!35!black,title={Before you move on}]
What does \zh{起因} mean? (\textquotedblleft{}What set it off; the starting cause\textquotedblright{}.) Say it once more.
\end{tcolorbox}
\section[\texorpdfstring{\zh{原来} (yuánlái)}{原来 (yuánlái)}]{\zh{原来} (yuánlái) — so (it turns out); as it turns out}

\label{lesson:ZH-C537-yuanlai}



\begin{quote}
Before the new one: say the Mandarin for something to look forward to, then the Mandarin for what set it off.
\end{quote}

\subsection*{You'll want to know: \zh{原来}}
\textbf{\zh{原来}} — \emph{yuánlái} — \textquotedblleft{}so (it turns out); as it turns out\textquotedblright{}.

\begin{quote}
\textbf{\zh{原来你也在这儿}} \emph{yuánlái nǐ yě zài zhèr} — So you're here too!
\end{quote}

The new character, in pieces: \zh{原} — ten strokes: the cliff, a bar and a long stroke falling to the left, as at the top of \zh{愿}; under it \zh{白}, five strokes, and \zh{小}, three strokes: \zh{愿} without its heart. The next lesson puts it in your hand.

\begin{cousinweb}
\textbf{\zh{原来}} (\emph{yuánlái}) is built from \zh{原} and \zh{来} 'to come', both inherited from Middle Chinese, \emph{ngjwon} and \emph{loj} (Guangyun, in Baxter's transcription). The dictionary glosses the first as wide flat land, a plain; the sense 'origin' is not given there, and how the pair came to mean 'as it turns out' is not settled.
\end{cousinweb}

\begin{grammarlens}[title={What you've built}]
One more word for saying what would happen.
\end{grammarlens}

\subsection*{Your turn}
\begin{itemize}
\raggedright
  \item \emph{Say it:} \emph{yuánlái}
  \item \emph{Say it:} \emph{yuánlái}, once more
  \item \emph{Say it:} \emph{yuánlái nǐ yě zài zhèr} — So you're here too!
  \item \emph{Recall:} say the Mandarin for something to look forward to, then the Mandarin for what set it off, then say \emph{yuánlái} again
\end{itemize}

\begin{tcolorbox}[breakable,colback=teal!4,colframe=teal!35!black,title={Before you move on}]
What does \zh{原来} mean? (\textquotedblleft{}So (it turns out); as it turns out\textquotedblright{}.) Say it once more.
\end{tcolorbox}
\section[\texorpdfstring{\zh{原} (yuán)}{原 (yuán)}]{\zh{原} — the cliff over \zh{白} and \zh{小}}

\label{lesson:ZH-W537-yuan-origin}



\begin{quote}
The word first, then the character inside it.

\begin{itemize}
\raggedright
  \item \emph{Recall:} say \emph{yuánlái} — \textbf{R1}, one lesson back, before it decays
\end{itemize}

\begin{quote}
\zh{愿} \zh{白}
\end{quote}

\emph{Read it:} each one aloud, with nothing said first — \emph{yuàn}, \emph{bái}
\end{quote}

\subsection*{Script — the cliff over \zh{白} and \zh{小}}
\hlblockfigure{\detokenize{figures/ZH-W537-yuan-origin-filmstrip.pdf}}{How \zh{原} is written, stroke by stroke}

\begin{quote}
\zh{原}
\end{quote}

\textbf{Three pieces: the cliff, then \zh{白}, then \zh{小} under it.} Count the pieces, not the strokes.

\textbf{The cliff, two strokes}: a bar, then a long stroke falling to the left down from its left end.

\textbf{\zh{白}, five strokes}, under the cliff's arm: a short stroke falling to the left, a vertical, a bar that turns down, two short bars.

\textbf{\zh{小}, three strokes}: a vertical that hooks to the left, a short stroke on its left, a dot on its right.

2 + 5 + 3 = ten, the first ten strokes of \zh{愿}.

\subsection*{Your turn}
\begin{itemize}
\raggedright
  \item \emph{Write it:} \zh{原}, all ten, with the model in front of you
  \item \emph{Cover:} write it again from memory, and count the pieces, not the strokes
\end{itemize}

\begin{tcolorbox}[breakable,colback=teal!4,colframe=teal!35!black,title={Before you move on}]
Which character is \zh{原} with \zh{心} under it? (\textbf{\zh{愿}}.) How many strokes does \zh{原} have? (\textbf{Ten}.)

Source: \href{https://raw.githubusercontent.com/chanind/hanzi-writer-data/68d10a4b21150cae5e1ebbd223eed289cf32d90c/data/\%E5\%8E\%9F.json}{Hanzi Writer Data \zh{原}, pinned stroke paths and medians}.
\end{tcolorbox}
\section[dǎozhì]{dǎozhì — to result in something bad}

\label{lesson:ZH-C537-daozhi}



\begin{quote}
Before the new one: say the Mandarin for what set it off, then the Mandarin for as it turns out.
\end{quote}

\subsection*{You'll want to know: dǎozhì}
\emph{dǎozhì} — \textquotedblleft{}to result in something bad\textquotedblright{}.

\begin{quote}
\emph{rúguǒ bú zhùyì jiù huì dǎozhì shìgù} — If you're careless, it will lead to an accident.
\end{quote}

Its characters come later; for now it is a word for the ear.

\begin{cousinweb}
\emph{dǎozhì} is built from \emph{dǎo} 'to lead' and \emph{zhì} 'to bring to', both inherited from Middle Chinese, \emph{dawH} and \emph{trijH} (Guangyun, in Baxter's transcription), which glosses the first 'to lead' and the second 'to arrive' and 'to send to'. The dipping tone of dǎo is not the regular outcome of \emph{dawH}.
\end{cousinweb}

\begin{grammarlens}[title={What you've built}]
One more word for saying what would happen.
\end{grammarlens}

\subsection*{Your turn}
\begin{itemize}
\raggedright
  \item \emph{Say it:} \emph{dǎozhì}
  \item \emph{Say it:} \emph{dǎozhì}, once more
  \item \emph{Say it:} \emph{rúguǒ bú zhùyì jiù huì dǎozhì shìgù} — If you're careless, it will lead to an accident.
  \item \emph{Recall:} say the Mandarin for what set it off, then the Mandarin for as it turns out, then say \emph{dǎozhì} again
\end{itemize}

\begin{tcolorbox}[breakable,colback=teal!4,colframe=teal!35!black,title={Before you move on}]
What does dǎozhì mean? (\textquotedblleft{}To result in something bad\textquotedblright{}.) Say it once more.
\end{tcolorbox}
\section[zàochéng]{zàochéng — to cause damage or trouble}

\label{lesson:ZH-C537-zaocheng}



\begin{quote}
Before the new one: say the Mandarin for as it turns out, then the Mandarin for to result in something bad.
\end{quote}

\subsection*{You'll want to know: zàochéng}
\emph{zàochéng} — \textquotedblleft{}to cause damage or trouble\textquotedblright{}.

\begin{quote}
\emph{rúguǒ nǐ bù shuō jiù huì zàochéng wùhuì} — If you don't say anything, it will cause a misunderstanding.
\end{quote}

Its characters come later; for now it is a word for the ear.

\begin{cousinweb}
\emph{zàochéng} is built from \emph{zào} 'to make' and \zh{成} 'to complete', both inherited from Middle Chinese, \emph{dzawX} and \emph{dzyeng} (Guangyun, in Baxter's transcription), which glosses the first 'to make'.
\end{cousinweb}

\begin{grammarlens}[title={What you've built}]
One more word for saying what would happen.
\end{grammarlens}

\subsection*{Your turn}
\begin{itemize}
\raggedright
  \item \emph{Say it:} \emph{zàochéng}
  \item \emph{Say it:} \emph{zàochéng}, once more
  \item \emph{Say it:} \emph{rúguǒ nǐ bù shuō jiù huì zàochéng wùhuì} — If you don't say anything, it will cause a misunderstanding.
  \item \emph{Recall:} say the Mandarin for as it turns out, then the Mandarin for to result in something bad, then say \emph{zàochéng} again
\end{itemize}

\begin{tcolorbox}[breakable,colback=teal!4,colframe=teal!35!black,title={Before you move on}]
What does zàochéng mean? (\textquotedblleft{}To cause damage or trouble\textquotedblright{}.) Say it once more.
\end{tcolorbox}
\section[xiàoguǒ]{xiàoguǒ — the effect it has (how well it works)}

\label{lesson:ZH-C537-xiaoguo}



\begin{quote}
Before the new one: say the Mandarin for to result in something bad, then the Mandarin for to cause damage.
\end{quote}

\subsection*{You'll want to know: xiàoguǒ}
\emph{xiàoguǒ} — \textquotedblleft{}the effect it has (how well it works)\textquotedblright{}.

\begin{quote}
\emph{yàoshi méiyǒu xiàoguǒ nǐ jiù huàn yí ge bànfǎ} — If it doesn't work, try another way.
\end{quote}

Its characters come later; for now it is a word for the ear.

\begin{cousinweb}
\emph{xiàoguǒ} is built from \emph{xiào} and \zh{果} 'fruit', both inherited from Middle Chinese, \emph{haewH} and \emph{kwaX} (Guangyun, in Baxter's transcription), which glosses the first 'to imitate' and 'proven effect': what the effort bears.
\end{cousinweb}

\begin{grammarlens}[title={What you've built}]
One more word for saying what would happen.
\end{grammarlens}

\subsection*{Your turn}
\begin{itemize}
\raggedright
  \item \emph{Say it:} \emph{xiàoguǒ}
  \item \emph{Say it:} \emph{xiàoguǒ}, once more
  \item \emph{Say it:} \emph{yàoshi méiyǒu xiàoguǒ nǐ jiù huàn yí ge bànfǎ} — If it doesn't work, try another way.
  \item \emph{Recall:} say the Mandarin for to result in something bad, then the Mandarin for to cause damage, then say \emph{xiàoguǒ} again
\end{itemize}

\begin{tcolorbox}[breakable,colback=teal!4,colframe=teal!35!black,title={Before you move on}]
What does xiàoguǒ mean? (\textquotedblleft{}The effect it has (how well it works)\textquotedblright{}.) Say it once more.
\end{tcolorbox}
